\documentclass[11pt,a4paper]{article}

\usepackage[utf8]{inputenc}
\usepackage[T1]{fontenc}
\usepackage{lmodern}
\usepackage{amsmath,amssymb,amsthm}
\usepackage{mathtools}
\usepackage[margin=1.1in]{geometry}
\usepackage{booktabs}
\usepackage{array}
\usepackage[hyphens]{url}
\usepackage{hyperref}

\theoremstyle{plain}
\newtheorem{theorem}{Theorem}
\newtheorem{proposition}[theorem]{Proposition}
\newtheorem{lemma}[theorem]{Lemma}
\theoremstyle{definition}
\newtheorem{definition}[theorem]{Definition}
\theoremstyle{remark}
\newtheorem{remark}[theorem]{Remark}

\newcommand{\FF}{\mathbb{F}}
\newcommand{\Sfe}{S_{\mathrm{fe}}}
\newcommand{\Dff}{D_{\mathrm{ff}}}
\newcommand{\Dfake}{D'_{\mathrm{ff}}}
\newcommand{\red}{\operatorname{red}}
\newcommand{\lean}[1]{\path{#1}}
\urlstyle{tt}
\emergencystretch=3em

\title{Machine-checked first-fall-degree reasoning\\ in Weil-descent ECDLP}
\author{Anonymous}
\date{31 October 2026}

\begin{document}
\maketitle

\begin{abstract}
Subexponential estimates for the elliptic-curve discrete logarithm problem over
small-characteristic fields rest on the claim that the polynomial systems from
Weil descent have a small \emph{first fall degree}. Several arguments in this
line replace the true first fall degree with a quantity that is easier to
compute, and the replacement is where errors enter. We give three results,
each checked in Lean~4 with Mathlib, and we state for each exactly what the
kernel checked and what it did not.
(i)~Lemma~4 of Nagao (ePrint 2015/984) asserts that the first fall degree is
at most a ``fake'' first fall degree computed modulo the field equations.
Under the equality reading of the paper's definition, we give two systems over
$\FF_2$, in two and three variables, with true first fall degree~3 and fake
first fall degree~2. A short argument shows when the inference can be
recovered for a given instance, and the paper's own applications satisfy it.
(ii)~We formalize Theorem~2.8 of Caminata and Gorla: if a system has a degree
fall and its last fall degree is finite, the last fall degree is the largest
degree at which a fall occurs. A single exhibited fall therefore gives a lower
bound and no upper bound.
(iii)~On a fixed Boolean Semaev $S_3$ system over $\FF_{2^4}$, the direct
Macaulay-row statistic and the degree-bounded closure agree on the first
degree with a new class but disagree on its multiplicity: 1~new class against
12 at degree three. Over 80 archived draws they agree on the grade of every
draw and disagree on the multiplicity for 77. The same audit withdraws an
earlier internal inference that every such system falls at degree two.
Nothing here gives an ECDLP algorithm or changes a published complexity
estimate.
\end{abstract}

\section{Introduction}

Semaev's summation polynomials~\cite{semaev04} turn a point decomposition on an
elliptic curve over $\FF_{q^n}$ into a polynomial system. Weil descent turns
that system into one over $\FF_q$ in more variables. Gr\"obner-basis
algorithms solve it in time governed by the degree they must reach, and the
\emph{first fall degree}, introduced in the HFE setting by Dubois and
Gama~\cite{duboisgama}, became a common proxy for that degree. Petit and
Quisquater~\cite{petitquisquater} conjectured that first fall degree tracks
solving degree for these systems. That conjecture, and its variants, underlie
the subexponential estimates of Semaev~\cite{semaev15} and
Nagao~\cite{nagao13,nagao15}. Huang, Kosters and Yeo~\cite{hky} and Kosters and
Yeo~\cite{kostersyeo} gave evidence against the assumption, and Caminata and
Gorla~\cite{cg} put the related invariants on a rigorous footing.

The arguments in this area pass between several quantities that differ
slightly: falls measured before or after reduction by the field equations;
field equations as generators of the system or as relations of the ring; row
spans of a Macaulay matrix or their multiplicative closure. Each step is easy
to get wrong and hard to check by eye. We used a proof assistant to check
three such steps. One check refutes a published lemma under a stated reading.
The second formalizes a known theorem in a form that makes the direction of
valid inference explicit. The third shows that two statistics with the same
name agree on one number and disagree on another.

\paragraph{What the kernel checked.} Every Lean statement cited here builds
against a pinned Mathlib and depends only on the axioms \texttt{propext},
\texttt{Classical.choice} and \texttt{Quot.sound}. No \texttt{sorry},
\texttt{admit}, custom axiom or \texttt{unsafe} declaration appears. Section~\ref{sec:status}
gives each statement's hosted CI receipt and the status of its independent
statement review, which is a separate check that the Lean proposition says
what the prose claims. The last-fall theorem, the concrete Semaev witness and
the trace identity of Section~\ref{sec:semaev} have passed that review at their
stated scope. The counterexample certificate and the Semaev rank family have
not.

\paragraph{What this paper does not claim.} We say nothing about Nagao's
Proposition~5 or about Semaev's first-fall assumption in general. We give no
complexity estimate and no algorithm, and no measurement here bears on
ECDLP cost at any cryptographic size.

\section{Definitions}\label{sec:defs}

Let $k$ be a field and $R=k[X_1,\dots,X_N]$. Write $\deg$ for total degree.

\begin{definition}[First fall degree, equality reading]\label{def:ff}
Let $S=\{s_1,\dots,s_M\}\subset R$. A \emph{fall at degree $D$} is a tuple
$(g_1,\dots,g_M)$ of polynomials with
\begin{enumerate}
\item $\deg(g_i s_i)\le D$ for every $i$, with equality for at least one $i$;
\item $\sum_i g_i s_i\neq 0$;
\item $\deg\bigl(\sum_i g_i s_i\bigr)<D$.
\end{enumerate}
The \emph{first fall degree} $\Dff(S)$ is the least $D\ge 1$ at which a fall
exists.
\end{definition}

Over $\FF_p$ let $\Sfe=\{X_i^p-X_i\}$ and write $\red(h)$ for the reduction of
$h$ modulo $\Sfe$, which over $\FF_2$ clamps every exponent to at most one.

\begin{definition}[Fake first fall degree]\label{def:fake}
$\Dfake(S)$ is defined as $\Dff(S)$ except that cofactors are reduced
(multilinear over $\FF_2$), every product and the sum are replaced by their
reductions, generators that reduce to zero are dropped, and condition~(2)
asks for a nonzero reduced sum.
\end{definition}

\begin{remark}[Which reading]\label{rem:reading}
Definitions~5 and~6 of~\cite{nagao15} print condition~(1) as
$\max_i\deg(g_is_i)\ge D$. Under that literal reading, padding the cofactors
makes condition~(1) free. The first fall degree then collapses to one more
than the least degree of a nonzero element of the ideal, and Lemma~4 holds
for reasons unrelated to the paper's argument. Nagao's earlier paper~\cite{nagao13} prints
equality in its Definition~1, and the standard definition used
in~\cite{cg} is the degree at which the fall happens. We adopt the equality
reading throughout. It is a declared interpretation of an ambiguous text. We
make no claim about the author's intent.
\end{remark}

\begin{definition}[Caminata--Gorla~\cite{cg}, Definition~1.5]\label{def:V}
For $F\subset R$ and $i\in\mathbb N$, let $V_{F,i}$ be the smallest
$k$-subspace of $R$ that contains every $f\in F$ with $\deg f\le i$ and is
closed under $f\mapsto fg$ whenever $f\in V_{F,i}$ and
$\deg f+\deg g\le i$. For $f$ in the ideal $(F)$ let
$d_f=\min\{d: f\in V_{F,d}\}$, so $d_f\ge\deg f$. Say $f$ \emph{has a degree
fall} if $d_f>\deg f$. The \emph{last fall degree} is
\[
d_F=\min\{d:\ f\in V_{F,\max(d,\deg f)}\ \text{for every } f\in(F)\}.
\]
\end{definition}

The Lean formalization writes the closure condition as
$\deg f+\deg g\le i$ rather than $\deg g\le i-\deg f$. On natural numbers the
second form truncates and would admit degree-zero multipliers of any $f$ with
$\deg f>i$.

\section{A counterexample to Lemma 4 of ePrint 2015/984}\label{sec:nagao}

Lemma~4 of~\cite{nagao15} states $\Dff\le\Dfake$, so that a bound on the easily
computed fake quantity would bound the true one. Two placements of $\Sfe$ are
possible on the true side. The printed lemma takes $\Dff$ of
$\{f_1,\dots,f_M\}$ alone. The \emph{inside form} takes both quantities of
the same system $S=F\cup\Sfe$, which is the form the paper's arguments
need. Lean checks the refutation of the inside form. The printed form fails
on the same witnesses: without $\Sfe$, $X^2Y+X(XY+X)=X^2$ and
$X_3(X_1X_2+X_3)+X_2(X_1X_3+X_2)=X_2^2+X_3^2$ are falls at degree~3, and the
degree-2 argument below applies unchanged. That check is by hand and by
exhaustive search, not in Lean.

\begin{theorem}\label{thm:nagao}
Over $\FF_2$, with $\Sfe$ among the generators of the true system, let
\begin{align*}
\text{(A)}\quad & F_A=\{X^2Y,\ XY+X\}\subset\FF_2[X,Y],\\
\text{(B)}\quad & F_B=\{X_1X_2+X_3,\ X_1X_3+X_2\}\subset\FF_2[X_1,X_2,X_3].
\end{align*}
For $S=F_A\cup\Sfe$ and for $S=F_B\cup\Sfe$ we have $\Dff(S)=3$ and
$\Dfake(S)=2$. Hence $\Dff\le\Dfake$ fails, and the universal statement over
finite systems over $\FF_2$ is false.
\end{theorem}

\begin{proof}
\emph{Fake side.} For (A), $\red(X^2Y)=XY$, and
$\red(1\cdot X^2Y)+\red(1\cdot(XY+X))=X$ has degree $1<2$, so there is a fake
fall at degree~2. For (B),
$\red\bigl((X_1+1)(X_1X_3+X_2)\bigr)=X_1X_2+X_2$, and adding $X_1X_2+X_3$ gives
$X_2+X_3$, again a fake fall at degree~2. No fake fall exists at degree~1,
because every reduced generator has degree~2.

\emph{True side, degree 3.} For (A), $X^2Y+X\cdot(XY+X)=X^2$. For (B),
$X_3\,(X_1X_3+X_2)+X_1\,(X_3^2+X_3)=X_2X_3+X_1X_3$. Each is a fall at degree~3.

\emph{True side, no fall below 3.} At $D=1$ every generator has degree at
least~2, so no product has degree exactly~1. At $D=2$ a product of degree at
most~2 needs a constant cofactor on a quadratic generator and the zero
cofactor on a cubic one. A fall then needs the degree-2 parts to cancel.
For (A) those parts are $XY$, $X^2$, $Y^2$, and for (B) they are $X_1X_2$,
$X_1X_3$, $X_1^2$, $X_2^2$, $X_3^2$. Both sets are linearly independent, so
no fall exists.
\end{proof}

The Lean module \lean{CryptoResearch.Nagao.Lemma4Witness} checks all of these
facts by \texttt{decide +kernel} in a model where a polynomial over $\FF_2$ is
a duplicate-free list of exponent vectors and cofactors range over polynomials
whose monomials have degree at most $D$ (multilinear on the fake side). The
named theorem \lean{lemma4_inside_false} states
$\neg\,\forall n\,S,\ \mathrm{Lemma4Inside}\ n\ S$. A separate Python
implementation, written from the definitions without reference to the
first, reproduces both witnesses. It also finds the lemma violated on 94 of
400 random pairs of multilinear quadratics in three variables
(23.5\%)~\cite{artifacts}. We reproduced both witnesses a third time by
exhaustive search over cofactor tuples while preparing this paper.

\subsection{When the inference survives}

The failure has one cause. Reduction modulo $\Sfe$ can lower the degree of an
individual product $g_is_i$, so a fake fall need not come from any true fall
at the same degree. The paper's Lemma~3 controls the degree of the sum and
says nothing about the products whose maximum defines $\Dff$. The following
proposition makes the repair explicit.

\begin{proposition}[Instance-wise recovery]\label{prop:recovery}
Over $\FF_2$, let $S=F\cup\Sfe$ and let $(g_i)$ be a fake fall at degree $v$
for the generators in $F$. If every unreduced product satisfies
$\deg(g_if_i)\le v$, then $\Dff(S)\le v$.
\end{proposition}

\begin{proof}
Let $s=\sum_ig_if_i$ and $r=\red(s)$, so $r\neq0$ and $\deg r<v$. Reduce $s$
one monomial at a time, replacing $X_j^2m$ by $X_jm$. Each step adds
$m(X_j^2+X_j)$, whose degree is at most the degree of the monomial being
reduced, and degrees never increase along the way. So
$r=s+\sum_jh_j(X_j^2+X_j)$ with $\deg\bigl(h_j(X_j^2+X_j)\bigr)\le\deg s\le v$.
This is a combination of the generators of $S$ in which every product has
degree at most $v$ and the sum $r$ is nonzero of degree less than $v$. Some
reduced product $\red(g_if_i)$ has degree exactly $v$, and the unreduced
product has at least that degree, so by hypothesis it has degree exactly $v$.
Hence there is a fall at degree $v$, and $\Dff(S)\le v$.
\end{proof}

Proposition~\ref{prop:recovery} is not formalized in Lean; the proof above is
the only check. It needs two hypotheses that neither~\cite{nagao13}
nor~\cite{nagao15} states as such: the field equations are generators of the
true system, and the fake witness loses no degree under reduction. In both
witnesses of Theorem~\ref{thm:nagao} the second hypothesis fails: $X^2Y$ drops
from degree~3 to~2, and $(X_1+1)(X_1X_3+X_2)$ drops from~3 to~2.

According to our reading of the source, the only use of Lemma~4
in~\cite{nagao15} is in the proof of its Lemma~6. That proof puts $\Sfe$ on
the true side and tracks product degrees with explicit cofactors of degree at
most one on reduced generators. Both hypotheses of
Proposition~\ref{prop:recovery} can be checked there.
\cite{nagao13} does not state Lemma~4 and counts the field equations among
its $n+N$ equations. In our reading, the published uses survive. What fails
is the lemma as a general tool: a later argument that cites Lemma~4 for a
fake-degree bound is relying on a false statement.

\section{The last fall degree is the largest fall}\label{sec:lastfall}

\begin{theorem}[Caminata--Gorla~\cite{cg}, Theorem~2.8; formalized]\label{thm:lastfall}
Let $k$ be a field and $F\subset k[X_1,\dots,X_N]$. Suppose $d_F$ is finite
and some $f\in(F)$ has a degree fall. Then
\[
d_F=\max\{\,d_f:\ f\in(F),\ d_f>\deg f\,\}.
\]
\end{theorem}

The Lean statement
\begin{center}\lean{CryptoResearch.FallDegree.lastFallDegree_eq_max_degreeFall}\end{center}
is over Mathlib's \texttt{MvPolynomial} for an arbitrary variable type and
arbitrary field. It asserts that $d_F$ is the greatest element of the set of
fall degrees. Finiteness of $d_F$ is an explicit hypothesis: Caminata and
Gorla obtain it from the existence of a Gr\"obner basis, and we did not
reprove that. The hypothesis that a fall exists is also explicit. The
formalization does not construct a system satisfying both. The source review
works through $F=\{x^2,\ x^2+x\}$, whose last fall degree is~2, on paper only.
The same review noted a local wording slip (``$f\in F$'') in the published
proof and supplied a direct derivation; the theorem is unaffected.

The module also proves the lower-bound direction on its own: if some $f$ in
the ideal has $d_f>\deg f$, then $d_F\ge d_f$. This is the direction an
experiment can use. Observing one fall certifies that the last fall degree is
at least the degree of that fall. It certifies nothing in the other
direction. An upper bound needs a closure argument covering every element of
the ideal, which a finite computation of a Macaulay rank profile does not
provide unless it is proved to compute the spaces $V_{F,i}$ and some bound
rules out later falls. Lemma~4 of~\cite{nagao15} was an attempt to get an
upper bound cheaply, in the first-fall setting. Section~\ref{sec:nagao} shows
that the shortcut fails.

\section{Two first-fall statistics on a Boolean Semaev system}\label{sec:semaev}

The third result is about finite computation. It shows how two statistics,
both reasonably called ``first fall'', separate on a real system.

\paragraph{The system.} Take the curve $y^2+xy=x^3+x^2+1$ over
$\FF_{2^4}=\FF_2[t]/(t^4+t+1)$ and the third summation polynomial
$S_3=e_2^2+e_3+1$ in the elementary symmetric functions of three abscissae.
Fix the target abscissa $x_R=1$ and restrict the two free abscissae to an
$\FF_2$-subspace with encoded basis $[9,3,2,12]$. Weil descent yields four
Boolean equations in eight variables $u_0,\dots,u_7$, in
$\FF_2[u]/(u_i^2-u_i)$. The archived equations have SHA-256
\path{89ab8e66...ba6eb}~\cite{artifacts}.

\paragraph{Two statistics.} The \emph{direct-row} statistic at degree $D$ is
the number of new low-degree classes,
$\dim(\mathrm{rows}_D\cap R_{<D})-\dim\mathrm{rows}_{D-1}$, where
$\mathrm{rows}_D$ is the span of the multiplier-times-equation rows of degree
at most $D$. The \emph{closure} statistic replaces $\mathrm{rows}_D$ by the
least subspace containing the equations and closed under multiplication
within the degree budget, the Boolean-quotient analogue of $V_{F,D}$.

\begin{theorem}[finite certificates]\label{thm:semaev}
For this system the direct-row ranks at degrees $1,2,3$ are $1,11,53$, with
top-degree projection ranks $1,10,41$, so the direct-row new-class counts are
$0,0,1$. There is an explicit vector in $\mathrm{rows}_3$ of degree at most two
that lies outside $\mathrm{rows}_2$.
\end{theorem}

Lean checks these ranks and the witness. The witness is an XOR of selected
degree-three rows, whose coefficient sum at monomial masks~80 and~192 vanishes
on every degree-two row but equals one on the witness. The modules are
\lean{CryptoResearch.Semaev.BooleanParity*} and
\lean{CryptoResearch.Semaev.ConcreteFirstFallWitness}, and they also check
that the four transcribed equations match a direct evaluation of $S_3$ on all
256 Boolean assignments.

An independent code review then computed the closure statistic on the same
system and, separately, on all 80 archived draws~\cite{artifacts}. Three
fixed-point methods agree on the closed spaces, checked by joint rank.

\begin{center}
\begin{tabular}{@{}rccc@{}}
\toprule
$D$ & direct rank / top / low & closed rank / top / low & new classes (direct, closed)\\
\midrule
1 & 1 / 1 / 0 & 1 / 1 / 0 & --\\
2 & 11 / 10 / 1 & 11 / 10 / 1 & 0,\ 0\\
3 & 53 / 41 / 12 & 79 / 56 / 23 & 1,\ 12\\
\bottomrule
\end{tabular}
\end{center}

On this draw the two statistics agree that the first new class appears at
degree three and disagree on how many appear: 1 against~12. Across the 80
draws they agree on the first degree with a positive count for every draw
(68 at degree two, 12 at degree three) and disagree on the count at degree
three for 77 draws. At degrees one and two, seven direct-row spaces contain
rows outside the closure. The direct-row policy admits a product based on
its degree after cancellation, even when the multiplier exceeds the degree
budget of the closure definition.

\paragraph{A withdrawn inference.} An earlier internal note proved in Lean a
trace identity: for $x_R\neq0$ and $\mu=x_R^{-2}$, one has
$\mu^{1/2}+\mu x_R=0$ in $\FF_{2^n}$
(\lean{trace_witness_annihilates_degree_two_part}). It cited the identity for
a nontrivial first fall at degree two for every admissible parameter choice.
The identity is correct. The inference does not follow, because the theorem
contains no statement about membership in a lower row space, and the zero map
satisfies its premises. The draw above has $x_R\neq0$ and no new class at
degree two, which refutes the uniform claim in this finite model. The same
review found that the historical counting code combined Boolean monomials by
set union where coefficients over $\FF_2$ must cancel in pairs; correcting
this changes the archived split of first-fall grades over the 80 draws from
70/10 to 68/12. Kosters and Yeo~\cite{kostersyeo} describe degree-two first
falls as usual, which is consistent with 68 of 80.

\section{Verification status}\label{sec:status}

\begin{center}
\small
\begin{tabular}{@{}p{0.34\linewidth}p{0.22\linewidth}p{0.36\linewidth}@{}}
\toprule
Statement & Kernel & Independent statement review\\
\midrule
Thm.~\ref{thm:nagao} (\lean{lemma4_inside_false}) & hosted receipt, CI run 41155781140 & not yet reviewed; producer caps fidelity at partial (two scoping rows)\\
Prop.~\ref{prop:recovery} & not formalized & paper proof only\\
Thm.~\ref{thm:lastfall} & hosted receipt, CI run 41191503149 & validated at stated scope (three-joint review)\\
Thm.~\ref{thm:semaev}, concrete witness & hosted receipt, CI run 41155781140 & validated at scope, one draw\\
Thm.~\ref{thm:semaev}, rank family & hosted receipts, CI run 41155781140 & integrity audit only\\
Trace identity (\S\ref{sec:semaev}) & hosted receipt, CI run 41155781140 & identity valid; uniform inference withdrawn\\
Closure counts (\S\ref{sec:semaev}) & not formalized & independent code review\\
\bottomrule
\end{tabular}
\end{center}

A hosted receipt records that the workspace built, the per-theorem audit
found the named declaration with only the three standard axioms, and the
target specification and module hashes matched the frozen inventory. Each
receipt's own field for semantic review reads \emph{pending}; reviews are
filed separately and are summarized in the last column.

\section{Discussion}

The three results point one way. In first-fall arguments the dangerous
step is moving between a quantity that is easy to measure and the one
that matters: reduced against unreduced products, row spans against their
closure, one observed fall against a bound on all of them. A kernel catches a
false lemma only once someone writes the lemma down precisely enough to
check. An independent statement review catches a true theorem that is cited
for more than it says. We needed both: the first for Theorem~\ref{thm:nagao},
the second for the withdrawn trace inference.

Three items would strengthen this note. The first is a Lean proof of the
printed (outside) form of the counterexample and of
Proposition~\ref{prop:recovery}. The second is a bridge from the
list-of-monomials model to \texttt{MvPolynomial}. The third is an independent
review of the Lean statement of Theorem~\ref{thm:nagao}.

\begin{thebibliography}{99}

\bibitem{cg}
A.~Caminata and E.~Gorla.
\newblock Solving degree, last fall degree, and related invariants.
\newblock \emph{J. Symbolic Comput.}; IACR ePrint 2021/1611; arXiv:2112.05579.

\bibitem{duboisgama}
V.~Dubois and N.~Gama.
\newblock The degree of regularity of {HFE} systems.
\newblock In \emph{ASIACRYPT 2010}, LNCS 6477, 2010.

\bibitem{hky}
M.-D. Huang, M.~Kosters, and S.~L. Yeo.
\newblock Last fall degree, {HFE}, and {W}eil descent attacks on {ECDLP}.
\newblock In \emph{CRYPTO 2015}; IACR ePrint 2015/573.

\bibitem{kostersyeo}
M.~Kosters and S.~L. Yeo.
\newblock Notes on summation polynomials.
\newblock arXiv:1503.08001, 2015.

\bibitem{nagao13}
K.~Nagao.
\newblock Equations system coming from {W}eil descent and subexponential
  attack for algebraic curve cryptosystem.
\newblock IACR ePrint 2013/549.

\bibitem{nagao15}
K.~Nagao.
\newblock Complexity of {ECDLP} under the first fall degree assumption.
\newblock IACR ePrint 2015/984.

\bibitem{petitquisquater}
C.~Petit and J.-J. Quisquater.
\newblock On polynomial systems arising from a {W}eil descent.
\newblock In \emph{ASIACRYPT 2012}, LNCS 7658, 2012.

\bibitem{semaev04}
I.~Semaev.
\newblock Summation polynomials and the discrete logarithm problem on
  elliptic curves.
\newblock IACR ePrint 2004/031.

\bibitem{semaev15}
I.~Semaev.
\newblock New algorithm for the discrete logarithm problem on elliptic
  curves.
\newblock IACR ePrint 2015/310; arXiv:1504.01175.

\bibitem{artifacts}
Artifact repository for this note (anonymized for review). Lean sources
under \path{formal/CryptoResearch/}; counterexample certificate
\path{tools/lemma4_inside_form.py}; closure audit
\path{docs/reports/2026-10-26-semaev-closure-review/}; hosted receipts
\path{docs/reports/2026-10-31-hosted-formal-receipts/}.

\end{thebibliography}

\end{document}
